\begin{center}
\fontsize{16pt}{14pt}\selectfont
\bfseries Лопатин Александр Сергеевич\mdseries

\fontsize{14pt}{14pt}\selectfont
Резюме
\end{center}
{
\fontsize{12pt}{12pt}\selectfont
\begin{flushright}
\item дата рождения: 10 июня 1990
\item мобильный телефон: +7 911 83 66445
%\item e-mail: \href{mailto:sbar.geek@gmail.com}{sbar.geek@gmail.com}
\item e-mail: \includegraphics[scale=0.2]{email}
\item примеры кода: \url{http://github.com/sbar}
\end{flushright}
}

\setcounter{section}{0}
\section{Прикладное программирование}
\begin{itemize}
\item C++ (ООП, темплейты, основные паттерны проектирования,
фреймворки Qt и Marmalade (aka AirPlay), базовые знания OpenGL),
кроссплатформенная разработка (GNU/Linux, Windows, iOS, Android и др.),
использование средств отладки и анализа (gdb, valgrind, cppcheck),
опыт работы с Visual~Studio
\item небольшой опыт работы с Java (J2ME, Android)
\end{itemize}

\section{Разработка скриптов}
\begin{itemize}
\item Python (сетевые клиент-серверные и веб-приложения, парсинг, базы данных и
другое)
\item UNIX shell (bash)
\item элементарные знания других скриптовых языков (Lua, Perl, PHP и другие)
\end{itemize}

\section{Веб-разработка}
\begin{itemize}
\item немного опыта разработки веб-приложений (Python, Django)
и верстки (HTML~+~CSS, LaTeX)
\item SQL (базовые знания, разработка реляционных баз данных)
в реализациях SQLite, MySQL и PostgreSQL
\end{itemize}

\section{Опыт работы}
\begin{itemize}
\item с 2012 ``Младший инженер-программист'' в SPL (SmartPhoneLabs),
Великий Новгород

Разработка и портирование компьютерных игр для мобильных платформ
(iOS, Android, BlackBerry, Bada и др.)

Выпущенный проект --- игра Panda Blair (iOS, Android, BlackBerry PlayBook)

\item с 2011 ``Младший сотрудник службы качества'' в SPL (SmartPhoneLabs),
Великий Новгород

Ручное тестирование мобильных Java-приложений
(Symbian, Windows~Mobile, BlackBerry и др.),
автоматизация внутренней деятельности посредством скриптинга на UNIX shell
и Python
\end{itemize}

\section{Образование}
\begin{itemize}
%\item с 2011 Северо-Западный Государственный Технический Университет
%(Санкт-Петербург), специальность
%``Программное обеспечение вычислительной техники и автоматизированных систем'',
%заочное~обучение
\item 2005---2009 Многопрофильный политехнический колледж НовГУ, специальность
``Программное обеспечение вычислительной техники и автоматизированных систем''
\end{itemize}

\section{Другое}
Есть умение понимать и изменять чужой код.

Владею английским языком на среднем уровне (достаточно для чтения/написания
документации и делового общения); родной язык --- русский.

Являюсь уверенным пользователем и разработчиком Gentoo GNU/Linux.

Есть опыт использования систем контроля версиями (git в основном) и
багтрекеров (Jira, Bugzilla, Launchpad, Mantis и др.)

Быстрообучаемость новым технологиям.
Прошел военную службу в Вооруженных Силах.
Не~курю.
